\documentclass{article}

\usepackage[T1]{fontenc}

\usepackage{booktabs,graphicx,siunitx}

\begin{document}

\section{Hausman test}

\begin{tabular}{llll}
\toprule
 & Fixed effects & Random effects & Difference \\
\midrule
investment & 0.13 & 0.40 & $-$0.27 \\
 & (0.04) & (0.10) &  \\
 & [0.04, 0.22] & [0.20, 0.59] &  \\
education\_spend & 0.15 & $-$1.57 & 1.72 \\
 & (0.16) & (0.24) &  \\
 & [$-$0.18, 0.47] & [$-$2.04, $-$1.10] &  \\
urbanization & $-$0.03 & 0.11 & $-$0.14 \\
 & (0.02) & (0.05) &  \\
 & [$-$0.06, 0.01] & [0.00, 0.22] &  \\
\midrule
Num.Obs. & 100 &  &  \\
N entities & 10 &  &  \\
R\textsuperscript{2} & .44 & .93 &  \\
\multicolumn{4}{l}{Hausman $\chi$\textsuperscript{2}(3) = 572.56, p $<$.001} \\
\bottomrule
\end{tabular}

\includegraphics{figures/01_hausman-test/figure_coefficients.png}

Tests whether the random-effects assumption holds. A p below alpha means the estimates differ systematically $\rightarrow$ use fixed effects; a non-significant p means random effects is acceptable (and more efficient). Method: R plm package (Croissant \& Millo, 2008). The standard errors in parentheses are clustered by entity; the bracketed line is the 95\% CI. Your significance threshold ($\alpha$) is 0.05.

A Hausman test comparing the fixed- and random-effects estimates gave $\chi$\textsuperscript{2}(3)=572.56, p $<$ .001.

Statistical basis: Hausman specification test (Hausman, J. A. (1978). "Specification tests in econometrics." Econometrica, 46(6), 1251-1271.); Implementation (plm package) (Croissant, Y., Millo, G. (2008). "Panel Data Econometrics in R: The plm Package." Journal of Statistical Software, 27(2), 1-43.).

\end{document}